\documentclass[11pt]{article}
\usepackage[margin=2.5cm]{geometry}
\usepackage{amsmath}
\usepackage{graphicx}
\usepackage{booktabs}
\usepackage[colorlinks,citecolor=blue,urlcolor=blue]{hyperref}

\title{One adsorption model for the phytate--oleosin interaction:\\
       \large reconciling $^{31}$P NMR, $\zeta$-potential and calorimetry}
\author{}
\date{}

\begin{document}
\maketitle

\section{Methods}

\subsection{Monte Carlo simulations}

Coarse-grained Monte Carlo simulations resolved the phytate--oleosin interaction
at the molecular scale.  The N-terminal arm of \textit{Cannabis sativa} oleosin
(UniProt~A0A803P977, \texttt{MKSFLPESGPSAQQ}) was represented at one bead per
residue with the CALVADOS~3 model~\cite{Tesei2021,Cao2024}; fifty chains were
end-grafted to one wall of a $130\times130\times200$\,\AA$^{3}$ slit at a surface
density of $3\times10^{-3}$\,\AA$^{-2}$, forming a brush that stands in for the
oil-body surface~\cite{Tzen1992}.  A single phytate molecule
(\textit{myo}-inositol hexakisphosphate) was modelled as a rigid seven-bead unit:
one inositol centre and six phosphates.  Beads interacted through the
Ashbaugh--Hatch potential~\cite{Ashbaugh2008} and a screened-Coulomb term set by
the salt, with phytate--peptide cross-interactions made purely repulsive.  Both
the peptide and phytate titrated with pH in the semi-grand canonical
ensemble~\cite{SmithTriska1994}, exchanging protons with an implicit bath through
molecular swap moves.  The peptide has four titratable groups (Glu, Lys, and the
two termini); phytate is less standard---its twelve dissociable protons (two per
phosphate) give thirteen charge states from neutral H$_{12}$PA to PA$^{12-}$, with
p$K_\text{a}$ values spanning $\sim\!1$ to $12$~\cite{DeStefano2003,Grimsley2009}.
Titrating all twelve reproduces the strong pH-dependence of the phytate valence
$z$, which sets the lateral self-repulsion of the adsorbed layer
($\hat a_2\propto z^2$, below).  Faunus~\cite{Stenqvist2013} sampled
configurations at pH~4, 6, 7 and 10 and 5--150\,mM NaCl.

The simulation yielded five profiles against the distance $r$ from the grafted
wall.  Umbrella sampling along the inositol $z$-coordinate (40 windows) gave the
potential of mean force $w(r)$, referenced to zero in the bulk.  The laterally
averaged electric potential $\varphi(r)$ gave a simulated $\zeta$-potential as
$\varphi$ at the outer edge of the grafted charge.  The rotational mean-square
displacement gave the relative rotational mobility $D/D_0$ of an adsorbed phytate.
A shape analysis gave the wall alignment $S_\text{wall}$, the second-rank
orientational order parameter of the disk-shaped phytate relative to the wall
($S_\text{wall}=0$ isotropic, $-\tfrac12$ face-on with the disk parallel to the
wall, $+1$ edge-on).  The fifth profile, the mean phytate charge, fixed the
pH-dependent valence $z$.  Full parameters, move sets, and p$K_\text{a}$ tables are
in the accompanying working notes.

\subsection{The gvdW adsorption model}

The potential of mean force is the wall potential of classical adsorption theory,
so the phytate accumulated at the brush follows from the generalised van der Waals
(gvdW) treatment of a fluid against an attractive wall~\cite{Nordholm2018} without
further simulation.  The zero-coverage (Henry) constant is the Boltzmann-weighted
surface excess,
\begin{equation}\label{eq:KH}
K_\text{H} = \int_0^\infty \big(e^{-w(r)/k_BT}-1\big)\,dr,
\end{equation}
where $r$ is the distance from the grafted wall and $w$ the referenced PMF.  The
two-dimensional isotherm then relates the areal excess $\Gamma$ to the Henry
excess $\Gamma_\text{HA}=K_\text{H}\,n_B$ ($n_B$ the bulk phytate density) through
a hard-disk and a lateral term,
\begin{equation}\label{eq:isotherm}
\Gamma = \Gamma_\text{HA}(1-\Gamma d^2)
         \exp\!\Big(\frac{2\hat a_2\Gamma}{k_BT}-\frac{\Gamma d^2}{1-\Gamma d^2}\Big),
\qquad
\hat a_2 = -\tfrac12\!\int_d^\infty\! 2\pi r\,u(r)\,dr,
\end{equation}
with $d\approx11$\,\AA\ the phytate diameter and $\hat a_2$ the lateral
binding-energy coefficient of the adsorbed layer for the phytate--phytate pair
potential $u(r)$.  The valence $z(\text{pH})$ is read from the simulated phytate
charge, so the framework carries no fitted binding constant.  Because $\Gamma$
appears on both sides, Eq.~\ref{eq:isotherm} is implicit; in the coverage
$\theta=\Gamma d^2$ it reads $\theta = \Gamma_\text{HA} d^2\,(1-\theta)\,
\exp\!\big(\alpha\theta-\theta/(1-\theta)\big)$ with $\alpha=2\hat a_2/(k_BT d^2)$,
and we solve it for $\theta$ by fixed-point iteration at each pH and salt.

For a screened-Coulomb repulsion $u(r)/k_BT = z^2 l_B\,e^{-\kappa r}/r$ (Bjerrum
length $l_B$, Debye length $1/\kappa$), the mean-field integral gives
$\hat a_2/k_BT = -\pi z^2 l_B\,e^{-\kappa d}/\kappa$, so $\hat a_2\propto z^2$.  This
mean-field form overestimates the like-charge repulsion between adsorbed phytates.
We correct $\hat a_2$ following Forsman's correlation-corrected Poisson--Boltzmann
theory~\cite{Forsman2004}: below a correlation-hole radius $R_c$ the divergent
Coulomb repulsion is replaced by its tangent.  $R_c$ is the nearest-neighbour
spacing of the adsorbed phytates, taken as Forsman's smaller of the density and
surface-charge estimates,
\begin{equation}\label{eq:Rc}
R_c = \min\!\Big(1/\sqrt{\pi\Gamma},\ \sqrt{z/(\pi|\sigma_\text{net}|)}\Big),
\end{equation}
capped at $2/\kappa$ and floored at the hard core $d$.  The density spacing
$1/\sqrt{\pi\Gamma}$ always bounds the hole, so the correction softens the repulsion
but never past the physical inter-phytate spacing; it therefore does not drive the
charge reversal that an uncapped surface-charge runaway would.  The simulated PMF is
almost purely electrostatic, so we add one salt-robust short-range attraction of
depth $\varepsilon$ at contact for the non-electrostatic binding---specific
phosphate--Lys/Arg ion-pairing---that the coarse-grained model under-represents;
$\varepsilon$ is fixed by the NMR and shared by all three probes.  The adsorbed
layer is solved self-consistently, and the $\zeta$-potential is read at the shear
plane from a nonlinear Poisson--Boltzmann solve of the adsorbed charge on the
simulated potential.  The valence $z(r)$ is the MC-regulated value.  Because the
$\zeta$ sample is diluted a thousandfold, its coverage is the isotherm
re-equilibrated at $[\text{PA}]/1000$.  The shear plane is the MC fixed-charge edge
of the brush (the $2\%$-of-peak rule of the potential profile), so no parameter is
fitted to the $\zeta$ or calorimetric data beyond the NMR-fixed $\varepsilon$.

\subsection{NMR analysis}

The $^{31}$P spectra show one population-averaged resonance, and the shallow PMF
well ($\sim\!2.6\,k_BT$ at $I=0.15$\,M) gives sub-millisecond residence
times---far shorter than the relaxation decay---so the rates are in fast exchange:
\begin{equation}\label{eq:twosite}
\Delta R_\text{obs} = f_\text{bound}\,\Delta R_\text{bound}
                    + (1-f_\text{bound})\,\Delta R_\text{free},
\end{equation}
where $\Delta R = R_{1\rho}-R_1$ isolates the slow reorientation.  Two rates are
measured at each pH---the mixture rate $\Delta R_\text{obs}$ of the oil-body sample
and the free-phytate rate $\Delta R_\text{free}$ of the control---and the
simulation supplies $f_\text{bound}$ from the coverage, $f_\text{bound}=\Gamma
a/n_B$, with the oil-body specific area $a=6\phi_\text{OB}/d_{43}$.  A measured
rate thus becomes a parameter-free test of the simulated adsorption.  The
free-phytate hydrodynamics enter here too: rotational and translational diffusion
coefficients, estimated for each protonation state from RDKit conformer
ensembles~\cite{RDKit} with Rotne--Prager--Yamakawa bead
hydrodynamics~\cite{Zuk2018,Waszkiewicz2023}, give the fast free-tumbling reference
against which $\Delta R_\text{obs}$ is read; the same hydrodynamic radius sets the
excluded diameter $d\approx11$\,\AA\ used in the isotherm (Eq.~\ref{eq:isotherm}).

\section{Results}

One coarse-grained adsorption---taken from the simulated PMF and carried to finite
concentration by the gvdW model---accounts for all three
measurements of phytate binding to the oil body: the $^{31}$P relaxation, the
$\zeta$-potential, and the calorimetry.  A single non-electrostatic binding energy,
fixed by the NMR, is shared across the three; everything else is read from the
simulation.

\subsection{The simulation}

The simulation sets the stage (Fig.~\ref{fig:mc}).  At low salt the PMF has a deep
attractive well for pH~4 and~6 ($-17$ and $-15\,k_BT$ near $r\approx27$\,\AA) and
turns repulsive at pH~10 (Fig.~\ref{fig:mc}A).  Read at the outer edge of the fixed
brush charge ($\approx\!40$\,\AA), the simulated $\zeta$-potential reproduces the
measured bare-oil-body values at pH~4 and~10 ($+40$ vs.\ $+38$\,mV; $-72$ vs.\
$-56$\,mV) but not at pH~6 ($+5$ vs.\ $-12$\,mV, Fig.~\ref{fig:mc}C): the
single-peptide model carries the cationic oleosin patch but not the anionic
phospholipid, so it misses the sign of the near-isoelectric surface.  The potential
rises higher still inside the brush---to $+52$\,mV at the pH-4 binding
minimum---and that stronger field is what drives the adsorption there.  An adsorbed
phytate nonetheless stays mobile, tumbling at about half its bulk rate
($D/D_0\approx0.5$--$0.7$, Fig.~\ref{fig:mc}B) with only a weak face-on bias
($S_\text{wall}\approx-0.005$).

\begin{figure}[h]
\centering
\includegraphics[width=0.92\textwidth]{../mc/composite_lowsalt.png}
\caption{Simulated phytate at the oleosin brush, low salt: PMF $w(r)$, phytate
  charge and potential profiles, wall alignment $S_\text{wall}$, and relative
  rotational mobility $D/D_0$ against distance from the grafted wall, by pH.}
\label{fig:mc}
\end{figure}

\subsection{The simulations account for the \texorpdfstring{$^{31}$P}{31P} relaxation}

Combining the isotherm with the relaxation tests the adsorption without a fitted
binding constant (Table~\ref{tab:gvdw}).  The binding signature is the difference
$\Delta R_\text{obs}-\Delta R_\text{free}$, strongest at pH~4.  Henry's law alone
gets the pH-4/pH-6 order backwards: $K_\text{H}$ rises from pH~4 to~6, predicting
\emph{more} binding at pH~6.  The lateral repulsion resolves it---scaling as $z^2$
(Eq.~\ref{eq:isotherm}), it crowds the more deprotonated pH-6 layer more strongly and
caps its coverage below pH~4.  The gvdW bound fraction $f_\text{gvdW}$ then falls
monotonically from pH~4 to~10 and matches the NMR-required $f_\text{req}$ in both
order and magnitude ($1.3$, $1.1$, $0.1\%$ against $1.3$, $1.0$, $0.2\%$), once the
non-electrostatic contact term is folded into $K_\text{H}$.  At pH~10 the strongly
negative brush all but excludes phytate, and the small residual binding matches the
weak measured effect.  Cross-polarisation gave no signal at any pH, consistent with
the retained mobility of the simulated adsorbed state: binding broadens the line by
co-tumbling with the micron-scale droplet, not by freezing the molecule.

\begin{table}[h]
\centering
\begin{tabular}{@{}c ccccc c ccc@{}}
\toprule
 & \multicolumn{5}{c}{Simulation (gvdW, $I=0.15$\,M)} & &
   \multicolumn{3}{c}{$^{31}$P relaxation} \\
\cmidrule(lr){2-6}\cmidrule(lr){8-10}
pH & $z$ & $K_\text{H}$/\AA & $\alpha$ & $f_\text{H}$/\% & $f_\text{gvdW}$/\% & &
 $\Delta R_\text{free}$ & $\Delta R_\text{obs}$ & $f_\text{req}$/\% \\
\midrule
4  & $-5.9$  & $925$  & $-27$ & $8.5$  & $1.31$ & & $162$ & $559$ & $1.33$ \\
6  & $-7.1$  & $1209$ & $-38$ & $10.8$ & $1.14$ & & $130$ & $430$ & $1.00$ \\
10 & $-10.5$ & $15$   & $-85$ & $0.15$ & $0.10$ & & $5$   & $52$  & $0.16$ \\
\bottomrule
\end{tabular}
\caption{Simulation and NMR quantities in the two-site analysis
  (Eq.~\ref{eq:twosite}) for the $50\%$ oil-body sample: phytate valence $z$,
  Henry constant $K_\text{H}$ (Eq.~\ref{eq:KH}), lateral coefficient $\alpha=2\hat
  a_2/(k_BT d^2)$, the Henry bound fraction $f_\text{H}=aK_\text{H}$, and the gvdW
  bound fraction $f_\text{gvdW}$ (Eq.~\ref{eq:isotherm}).  $K_\text{H}$ includes the
  non-electrostatic contact term and so is weakly attractive even at pH~10.  Rates
  $\Delta R_\text{free}$ (control) and $\Delta R_\text{obs}$ (oil body) in s$^{-1}$;
  $f_\text{req}$ is the bound fraction the two-site model requires at $\Delta
  R_\text{bound}=3\times10^{4}$\,s$^{-1}$.}
\label{tab:gvdw}
\end{table}

\subsection{The \texorpdfstring{$\zeta$}{zeta}-potential fit}

The gvdW model is compared to the measured $\zeta$-potential over
pH, phytate concentration and added salt (Fig.~\ref{fig:zeta}) with \emph{no
parameter fitted to the $\zeta$ data}: the coverage, the potential shift and the
shear plane all come from the simulation, and the single binding energy---the
short-range non-electrostatic contact term---is fixed by the NMR.  The shear plane
is the fixed-charge edge of the brush plus adsorbed phytate (the $2\%$-of-peak rule
of Fig.~\ref{fig:mc}C), at $42$, $40$ and $36$\,\AA\ for pH~4, 6 and~10---the brush
periphery, where the hydrodynamic slip of a polymer-coated particle sits.  Against
the 39 measured means the model gives an RMS of $7.5$\,mV; weighting each point by
its measurement error raises the residual to $15$\,mV, concentrated almost entirely
at pH~4.  The bare-oil-body $\zeta$ that anchors each series is itself reproduced by
the simulation: solving the same Poisson--Boltzmann relation for the MC peptide
charge recovers the measured bare $\zeta$ and its salt screening at pH~4
quantitatively, validating the brush electrostatics.  At pH~6 and~10 the
coarse-grained brush carries a near-isoelectric sign error and a high-salt
overscreening, so the bare $\zeta$ is taken from experiment and the model is tested
on the phytate-induced shift alone.

The agreement is strongly condition-dependent.  At pH~6 the model tracks both the
phytate and the salt series within the experimental scatter, and at pH~10 it
reproduces the salt series.  The pH-4 series carries by far the smallest error
bars---$\zeta$ near the isoelectric point is measured to $\pm0.1$--$0.3$\,mV---and it
is where the model fails: it holds at $+16$\,mV at $7.6$\,mM~PA while the data cross
zero to $-4$\,mV.  The pH-4 charge reversal is not reproduced, and this is the honest
edge of the faithful treatment: over-charging would require softening the
phytate--phytate correlations over a range beyond their nearest-neighbour spacing,
which the density-limited correlation correction forbids.  A weak pH-10 phytate
effect (${\sim}5$--$10$\,mV more negative) is likewise not captured.

\begin{figure}[h]
\centering
\includegraphics[width=0.92\textwidth]{../zeta/fig_zeta.png}
\caption{Measured $\zeta$-potential (mean markers over a $\pm$SD band, coloured by pH:
  blue pH~4, orange pH~6, green pH~10) with the gvdW model (open
  triangles, evaluated only at the measured conditions): (A) $\zeta$ vs.\ [PA] at low
  salt; (B) $\zeta$ vs.\ [NaCl] at $1.5$\,mM~PA.  The model adds the simulated phytate
  shift to the measured bare-OB $\zeta$ (its per-condition anchor), so it is drawn as
  points, not a line.  No parameter is fitted to $\zeta$---the shear plane is the MC
  fixed-charge edge of the brush and the binding is fixed by the NMR.  RMS $7.5$\,mV
  (error-weighted $15$\,mV) over 39 points: pH~6 and the salt series fall within
  experimental error, while the pH-4 charge reversal is not reproduced.}
\label{fig:zeta}
\end{figure}

\subsection{The calorimetry}

The gvdW model fits the pH-4 thermogram about as well as the
instrument's independent one-site model (Fig.~\ref{fig:itc}, Table~\ref{tab:itc}),
and it is the more appropriate picture: phytate binds a charged brush, not a set of
independent chemical sites.  The one-site stoichiometry $n=0.38$ is itself a fitted,
concentration-coupled quantity---between $0.32$ and $0.45$ across the three
replicates---not a measured count, and it carries no surface physics.  The gvdW model
instead \emph{fixes} the binding strength $K_\text{H}$ and the lateral repulsion
$\hat a_2$ from the simulated PMF---the same $\hat a_2$ that governs the
$\zeta$-potential and the NMR---so the binding is predicted rather than fitted,
leaving only the enthalpy $\Delta H$ and the peptide-to-oleosin area free.  That is
two free parameters against the one-site model's three, and the less correlated of
the two (Table~\ref{tab:itc}).

\begin{table}[h]
\centering
\small
\begin{tabular}{@{}lcc@{}}
\toprule
 & one-site & gvdW (MC) \\
\midrule
adsorption & independent sites & charged surface \\
binding & $K_a=(2.5\pm0.8)\times10^{5}$\,M$^{-1}$ & $K_\text{H}$ \emph{fixed} (PMF) \\
lateral $\hat a_2$ & absent & \emph{fixed} (PMF, Forsman) \\
capacity & $n=0.38\pm0.03$ & area $=(3.0\pm0.1)\times10^{3}\,$\AA$^2$ \\
$\Delta H$\,/\,kJ\,mol$^{-1}$ & $-47\pm5$ & $-34\pm1$ \\
\midrule
free parameters & 3 & \textbf{2} \\
RMS\,/\,\textmu J & $5.9$ & $7.4$ \\
\bottomrule
\end{tabular}
\caption{The instrument one-site model and the MC-grounded gvdW model fit to the same
  pH-4 thermogram (Fig.~\ref{fig:itc}) at the cell ionic strength ($\approx\!15$\,mM);
  fitted values are $\pm1$ standard error from the fit covariance.  The one-site model
  fills independent sites (stoichiometry $n$); the gvdW model adsorbs phytate onto the
  charged brush and \emph{fixes} the binding $K_\text{H}$ and the lateral repulsion
  $\hat a_2$ from the simulated PMF, leaving only the area and $\Delta H$ free.  The
  one-site parameters are strongly correlated ($K_a$--$\Delta H$ $0.88$,
  $\Delta H$--$n$ $0.85$), so their values are poorly determined; the two-parameter
  gvdW fit is the more constrained (area--$\Delta H$ $0.59$).}
\label{tab:itc}
\end{table}

The simulated binding transfers to the calorimetry only once the cell ionic
strength is set correctly.  The 5\% oil-body cell carries oleosin and phospholipid
counterions that raise the ionic strength to $\approx\!15$\,mM, far above the 3\,mM
buffer; since $K_\text{H}$ is exponentially salt-sensitive, the buffer value would
overstate the binding $\sim\!100$-fold, whereas at the cell value the simulated
isotherm reproduces the thermogram directly (RMS $6.7$\,\textmu J), and floating the
ionic strength returns $\approx\!18$\,mM.

\begin{figure}[h]
\centering
\includegraphics[width=0.7\textwidth]{../itc/fig_itc_gvdw.png}
\caption{pH-4 ITC thermogram (points) with the MC-grounded gvdW model (blue) and the
  instrument one-site model (grey).  The gvdW fit reads the binding from the simulated
  PMF and reproduces the thermogram with only $\Delta H$ and an area free.}
\label{fig:itc}
\end{figure}

One adsorption therefore underlies all three probes, each measured at its own ionic
strength---the calorimetry at $\approx\!15$\,mM, the $\zeta$-potential after
dilution, and the NMR at $0.15$\,M---with the model rooted in the simulated PMF and
carrying no parameter fitted to the electrokinetic or calorimetric data.

\bibliographystyle{unsrt}
\bibliography{references}
\end{document}
